% Supplementary material for "When the Interval Is Smaller Than the Instrument".
% NOT part of the main submission: this document accompanies the paper as supplementary
% material only. Each section is indexed in docs/supplement_index.md, which records what moved
% here from which revision of the main text, so nothing is lost.
\documentclass[manuscript,screen,review]{acmart}
\settopmatter{printacmref=false}
\renewcommand\footnotetextcopyrightpermission[1]{}
\newcommand{\kafka}{Kafka}
\newcommand{\redis}{Redis}
\newcommand{\tti}{TTI}
% Generated by scripts/emit_paper_numbers.py from the campaign ledger.
% Do not edit by hand -- run the script and commit the result.
% Every number below is recomputed from docs/results/external_campaigns_index.csv.
\newcommand{\ombRuns}{215}
\newcommand{\ombDiscarded}{11{,}207{,}668}
\newcommand{\ombKept}{12{,}155{,}820}
\newcommand{\ombNegatives}{41{,}403}
\newcommand{\ombKafkaDiscarded}{10{,}588{,}439}
\newcommand{\ombKafkaNegatives}{0}
\newcommand{\ombRedisDiscarded}{619{,}229}
\newcommand{\ombRedisNegatives}{41{,}403}
\newcommand{\ombRetentionMin}{0.00}
\newcommand{\ombRetentionMax}{100.00}

\begin{document}
\title{Supplementary Material: When the Interval Is Smaller Than the Instrument}
\author{Gustavo Pedro Ricou}
\affiliation{%
  \institution{School of Computer Science and Statistics, Trinity College Dublin}
  \city{Dublin}
  \country{Ireland}}
\email{pedrorig@tcd.ie}
\maketitle

This document is supplementary material and is not part of the main submission. Main-text
references of the form ``supplementary material S$n$'' resolve here. The provenance of every
section --- what moved here, from where, and when --- is indexed in the artefact at
\texttt{docs/supplement\_index.md}.

\section{S1: The replay-rate provenance gap, in full}
\label{supp:rateprovenance}

Auditing the corpus turned the same scrutiny on our own record-keeping, and it did not survive
it either. We report the finding here because the protocol in the main text's protocol section is
weaker without it.

Replay plans built from StatsBomb carry a \textbf{baked-in $120\times$ time compression}: a
plan's emission offset is already the match time divided by $120$. The producer's
\texttt{-{}-speedup} flag then multiplies that. So \texttt{-{}-speedup 1} against such a plan is
\emph{not} real time, it is $120\times$; true real time needs $1/120$. The same flag means
different things depending on which plan it is pointed at, and getting it wrong is silent: the
run completes, the wall time looks unremarkable, and the numbers stay plausible. We lost a
campaign to it before adding a pre-flight that derives the value from the plan and then checks
the achieved rate against elapsed wall time.

The part we cannot repair is retrospective. \textbf{No surviving artefact records the achieved
replay rate of any run we report.} Several superseded Testbed~A scenarios record the nominal
\texttt{-{}-speedup} \emph{flag}, which is not the same thing: converting a flag to a rate requires
the plan's compression, and every plan those files name has since been deleted. Everywhere else
the orchestrator
wrote \texttt{speedup} into a per-campaign summary that lived beside the raw per-run data, and
for the earliest corpora only the aggregated CSVs were kept. For the E1 corpus of the main text's powered-transport section the surviving evidence is contradictory: the
commit that introduced it states the runs were made at true real time, the campaign script
reconstructed afterwards specifies \texttt{-{}-speedup 10}, and the flag's meaning was corrected
$21$~hours later the same day. E1 was replayed at $1\times$, $120\times$ or $1200\times$, and we
cannot say which.

\paragraph{The rate is recoverable from the measurements, if not from the metadata.} Having
argued all paper that a physical constraint can decide what record-keeping cannot, we applied the
same move here, and it works. The start-up cost of the main text's attribution section is a clock: it
lasts a fixed $102.6$~ms of \emph{wall} time, so the number of events it swallows depends
directly on the replay rate. Counting how many events per run carry it therefore reads the rate
back out of the data.

The E1 runs sort into cells by how many events they matched, and one cell is diagnostic. Fifteen
\kafka{} runs matched exactly eight events, and their median scheduling lag is
$\mathbf{52.34}$~ms (range $51.60$--$53.01$) --- not $103$~ms, and not $1.6$~ms. A median of
eight values is the mean of the fourth and fifth. The only way to land midway between the two
modes is for exactly four of the eight to be late, so the fourth value is a fast event and the
fifth a slow one: $(1.59 + 103.5)/2 = 52.6$~ms against $52.34$ observed. Four is also exactly
what the loop trace shows at a verified real-time rate (the window table in the main text), and it is what
the plan predicts, since five events share $t_{\mathrm{sim}}{=}0$ and the first of them pays the
block rather than inheriting it.

That cell excludes the accelerated alternatives outright. At $120\times$ the prologue spans
$12.3$~s of match time and holds $16$ events; at $1200\times$, $123$~s and $112$ events. Under
either, all eight matched events would be inside it and the cell would read $103$~ms. It reads
half that. \textbf{E1 was replayed at true real time}, the commit record was right, and the
campaign script reconstructed afterwards describes a different, earlier campaign. The remaining
cells agree: runs matching five or seven events read $103$~ms, as four late events out of five or
seven must give.

\emph{The inference has since been checked against the original script, and it holds.} Archiving
the cloud driver before releasing it turned up the campaign as actually invoked, in a shell script
in the operator's home directory that no part of the repository had ever referenced. It specifies
\texttt{-{}-speedup 0.008333} --- $1/120$, against a plan already compressed $120\times$, which is
true real time --- and its log's first line, \texttt{E1 N=1 reps=8 00:39:40}, names the run
directory the E1 artefact's first row still points at. The arithmetic above and the flag agree.

We report that as a check rather than as the answer, because the order matters. The rate was
determined from the measurements while the metadata was genuinely unavailable, and the script
surfaced afterwards, by luck, from a machine that was about to be switched off. What the
agreement establishes is not that E1's provenance was fine all along --- it was not --- but that
the recovery method returns the right answer when there is an answer to check it against. That is
the one thing a physical-constraint argument can rarely demonstrate about itself.

The same script also explains a feature we had described and not accounted for. It passes
\texttt{-{}-max-t-sim 2}: two seconds of match time per run. The median of seven matched events
that makes E1 under-powered, and that the main text's attribution section treats as an opening burst,
is that window. The burst is real and is why those seven events are the densest in the match; the
reason there were only seven is that the campaign asked for two seconds.

We leave the disclosure standing rather than deleting it, for two reasons. The recovery was
available only because the artefact we were chasing happened to be a wall-clock constant with a
sharp signature, and the confirmation only because a virtual machine had not yet been destroyed;
a reader should not have to reconstruct an experimental parameter by arithmetic on its results,
and next time there may be nothing to reconstruct it from and nothing left to check it against.
And the exercise makes the protocol rule concrete: what was missing was one number, written down
once.

\paragraph{What the gap touched.} The audit is unaffected either way: a negative transport is
causally impossible at any replay rate, and the rejection counts are arithmetic on surviving
per-event data. The withdrawal in the main text's attribution section is unaffected by construction ---
we give the argument in a form that holds at every candidate rate. The rules of
the main text's rules section were measured after the pre-flight existed, at a derived and verified
rate. the main text's powered-transport section's external validity is restored by the recovery above, and the
recovered script now documents it as well --- though the inference stood on its own before the
script appeared, and is what we would still be relying on had the machine been released a day
earlier.
the main text's protocol section carries the rule that would have made all of this unnecessary.



\section{S2: The load-axis post-mortem, in full}
\label{supp:loadaxis}

Moved from the main text's mixture development: the load-axis post-mortem in full. The main text keeps the admissions and the numbers.

\paragraph{Both load-axis predictions failed, including ours.} The knee sweep placed conditions
at achieved $\rho$ of $0.881$, $0.920$, $0.950$, $0.970$ and $0.990$, where the candidate forms
diverge most; the earlier ladder stopped at $0.878$ and everything above it collapsed onto a
degenerate $\rho = 1.000$, which is why the main text's rules section could not settle the question.
With those points in hand, M/G/1 fits \emph{worse than the mean} ($R^2 = -0.05$ against a fitted
exponential's $0.93$). The withdrawal does not merely stand: the forms are now separated, and
M/G/1 is the one ruled out.

The bounded prediction fared better but also failed. Over $\rho = 0.881$ to $0.990$ the rate grew
by a factor of $1.44$. M/G/1 predicted $13.8$; the bracket we registered before the sweep
predicted $2.45$--$3.07$. The rate saturates, as a bounded mechanism requires --- but harder than
our own extrapolation allowed, and a prediction that misses low is still a prediction that
missed. We report it as a failure rather than as the nearer of two misses.

The two failures have one cause, and E-A5 names it. Our bracket was still parameterised in
$\rho$, through $p = \rho^C$; E-A5 shows $\rho$ does not determine occupancy. What survives is
the mechanism --- established by manipulation --- and not any curve in $\rho$, ours included.


\section{S3: The campaign ledger schema}
\label{supp:ledger}

One row per cell of every instrumented-benchmark campaign,
\texttt{docs/results/external\_campaigns\_index.csv}:

\begin{itemize}
\item \texttt{campaign}, \texttt{cell} --- sub-campaign directory and cell name; the cell name
  encodes the swept axis (\texttt{l} load percent, \texttt{s} message size, \texttt{r} producer
  rate) and replicate number, e.g.\ \texttt{r625\_rep3}.
\item \texttt{level} --- the numeric value of the swept axis.
\item \texttt{valid} --- 1 if the cell's counts are usable under the pre-stated rule; 0 rows are
  retained with a reason, never deleted.
\item \texttt{count\_source} --- \texttt{shutdown\_hook} for exact end-of-run counters;
  \texttt{periodic\_quantised} for in-run lines (quantised to $10^4$, excluded from analysis);
  \texttt{none\_in\_log} when no counters were written.
\item \texttt{kept}, \texttt{discarded\_zero}, \texttt{discarded\_negative} --- the instrumented
  guard's three counters, separated by sign.
\item \texttt{log\_bytes}, \texttt{log\_sha256} --- size and hash of the run log the row was
  built from, so any row can be re-derived and checked against its source.
\end{itemize}

The Python-harness campaigns (no JVM, Section on generality in the main text) use the same cell
naming and publish \texttt{docs/results/external/harness\_results.csv} with per-cell counts,
sign-separated discards, and pacer-jitter percentiles.

\section{S4: OpenMessaging distributed mode, failure diagnostics}
\label{supp:distdiag}

Eleven attempts in total never produced a measurement. The three diagnosed attempts of the
referee campaign (chain 18, block E) fail identically at OMB commit \texttt{5b1fa70}:
\texttt{java.util.concurrent.CompletionException} wrapping
\texttt{java.lang.IllegalArgumentException} in the coordinator--worker path, before any workload
starts. Per attempt, the version, full coordinator and worker logs, and the extracted signature
are archived in \texttt{docs/results/external/dist\_diag/}; a draft upstream report is at
\texttt{docs/omb\_distributed\_issue.md} (not filed --- filing is the author's action).

\section{S5: The E1 reconciliation, in full}
\label{supp:e1rec}

\begin{table}[t]
\caption{The same runs reproduce both published answers. Transport medians (ms) from the powered
campaign, computed over every event and over each run's first seven events in emission order ---
the window E1 saw. The all-events columns reproduce the powered $0.41$~ms shift; the prologue
columns reproduce E1's near-equality (the main text) in both shift and absolute level. The
window, not the brokers, accounts for the disagreement.}
\label{tab:e1rep}
\begin{tabular}{rrrrrrrr}
\toprule
& & \multicolumn{3}{c}{All events} & \multicolumn{3}{c}{First $7$ (prologue)} \\
\cmidrule(lr){3-5}\cmidrule(lr){6-8}
$N$ & Runs & \kafka{} & \redis{} & Shift & \kafka{} & \redis{} & Shift \\
\midrule
1  & 5  & 0.476 & 0.095 & 0.381 & 0.831 & 0.743 & 0.088 \\
9  & 45 & 0.533 & 0.116 & 0.417 & 0.939 & 0.810 & 0.129 \\
12 & 55 & 0.527 & 0.113 & 0.414 & 1.018 & 0.769 & 0.248 \\
\bottomrule
\end{tabular}
\end{table}

Moved from the main text's powered-transport section: the complete windowed re-analysis showing E1's near-equality is the observation window and nothing else.

\paragraph{Why E1 saw near-equality, tested directly.} The two campaigns disagree by more than
power: E1 puts both systems near $0.8$~ms and within $0.1$~ms of each other, while the powered
campaign puts them at $0.53$ and $0.11$~ms and $0.41$~ms apart. Attributing that to seven events
is an explanation only if it survives being checked, so we checked it. If the difference is the
observation window and nothing else, then restricting the \emph{powered} runs to their first seven
events --- in emission order, the same events E1 would have caught --- must reproduce E1, and the
same runs must give the powered answer over their full length. That is a prediction with one free
parameter fixed in advance, and it can fail.

It does not fail (Table~\ref{tab:e1rep}). Over all events the runs give shifts of $0.381$,
$0.417$ and $0.414$~ms, reproducing the powered $0.41$~ms. Over their first seven they give
$0.088$, $0.129$ and $0.248$~ms --- E1's wandering near-equality, and of the same magnitude as
E1's own $0.021$, $0.116$ and $0.053$~ms. The \emph{absolute} medians match as well, which the
shift alone does not force: the prologue puts \kafka{} at $0.83$--$1.02$~ms and \redis{} at
$0.74$--$0.81$~ms, against E1's $0.79$--$1.00$ and $0.72$--$0.86$. One set of runs, one rate, one
protocol, two windows, both published answers.

So E1 is not wrong and the campaigns do not contradict each other. E1 measured the opening burst,
during which both systems are slow for the reason the main text's attribution section establishes ---
topic auto-creation, connection setup, and the first blocking \texttt{produce()} --- and that
shared start-up cost swamps a $0.41$~ms difference. As the run lengthens the cost dilutes, and
the brokers' steady-state separation emerges. The seven-event corpus answered a question about
start-up while appearing to answer one about brokers, which is the same failure of window
selection as the main text's attribution section, one level up.




\end{document}
